\documentclass[11pt]{article}
\usepackage[margin=1in]{geometry}
\usepackage[T1]{fontenc}
\usepackage[utf8]{inputenc}
\usepackage{lmodern}
\usepackage{xcolor}
\usepackage{booktabs}
\usepackage{longtable}
\usepackage{array}
\usepackage{amsmath,amssymb}
\usepackage{enumitem}
\usepackage{hyperref}
\hypersetup{colorlinks=true, linkcolor=blue, urlcolor=blue}
\definecolor{accent}{RGB}{25,56,99}

\setlength{\parskip}{0.45em}
\setlength{\parindent}{0pt}

\begin{document}

\begin{center}
{\LARGE\bfseries\color{accent} Comprehensive Comparison of Power-Law and Logistic Models for Saturation--Electrical Property Relationships in SIP Experiments}\\[0.4em]
{\large Pooled analysis of sample-level, site-level, and campaign-level trends using updated SIP comparison workflows}
\end{center}

\section{Abstract}

Understanding how electrical properties vary with water saturation is essential for interpreting Spectral Induced Polarization (SIP) experiments in partially saturated porous media. This report compares two functional models for describing saturation--electrical property relationships:

\begin{itemize}[leftmargin=1.5em]
\item a power-law model
\item a four-parameter logistic model
\end{itemize}

Model performance is evaluated using coefficient of determination ($R^2$), sum of squared errors (SSE), Gaussian log-likelihood ($\ln L$), residual standard deviation ($\sigma$), and Akaike Information Criterion (AIC). The updated workflow supports sample-level, site-level, and all-data comparisons, optional natural-log transformations, and normalization performed on the relevant aggregation range only.

The overall pattern is that resistivity--saturation relationships often remain well described by power-law scaling, while conductivity--saturation relationships can show sigmoidal or threshold-driven behavior that is better captured by logistic models. This distinction is especially important when conductivity rises sharply across a limited saturation interval before approaching a plateau.

\section{Introduction}

Electrical geophysical measurements are widely used to characterize subsurface hydrological properties. Spectral Induced Polarization (SIP) experiments provide valuable insight into how electrical conductivity and resistivity vary with fluid saturation in porous materials.

Understanding these relationships is important for:

\begin{itemize}[leftmargin=1.5em]
\item interpreting laboratory SIP column experiments
\item improving petrophysical models
\item linking electrical behavior to pore-scale processes
\item distinguishing monotonic scaling from threshold-controlled transitions
\end{itemize}

Historically, resistivity--saturation relationships have often been described using power-law scaling similar to Archie's law. Conductivity, however, may exhibit nonlinear transitions associated with connectivity thresholds, wetting-film development, or rapid changes in conductive pathway continuity. A logistic model provides a compact way to represent such behavior.

The current comparison workflow extends earlier analyses by introducing:
\begin{itemize}[leftmargin=1.5em]
\item direct comparison of power-law and logistic models across multiple aggregation levels
\item Gaussian likelihood-based statistics in addition to $R^2$
\item optional base-10 or natural-log transformations
\item group-consistent normalization so that sample-, site-, and global fits are each normalized only over their own data range
\end{itemize}

\section{Dataset description}

The comparison framework uses SIP outputs previously generated by the main processing pipeline. The combined dataset may include measurements from multiple sites, such as:

\begin{itemize}[leftmargin=1.5em]
\item Rosemead
\item Van Nuys (VN)
\end{itemize}

Each joined dataset contains measurements of:
\begin{itemize}[leftmargin=1.5em]
\item water saturation ($S$)
\item electrical resistivity ($\rho$)
\item electrical conductivity ($\sigma$)
\item their transformed or normalized counterparts
\end{itemize}

The model-comparison scripts operate at three analysis levels:
\begin{longtable}{p{0.18\textwidth} p{0.68\textwidth}}
\toprule
\textbf{Level} & \textbf{Definition}\\
\midrule
\endhead
Sample & Each run folder is analyzed separately.\\
Site & All runs belonging to the same site are pooled together.\\
All & All available runs are combined into one global dataset.\\
\bottomrule
\end{longtable}

This hierarchy makes it possible to distinguish between specimen-specific behavior, site-specific tendencies, and broader campaign-wide trends.

\section{Model formulations}

\subsection{Power-law model}

In real space, the power-law model is written as:
\[
y = aS^{b}
\]

where:
\begin{itemize}[leftmargin=1.5em]
\item $S$ is water saturation
\item $a$ is a scaling coefficient
\item $b$ is the saturation exponent
\end{itemize}

When a logarithmic transformation is used, the same relationship becomes linear:
\[
\log y = b \log S + \log a
\]

The workflow supports either:
\[
\log_{10}(\cdot)
\qquad \text{or} \qquad
\ln(\cdot)
\]
depending on the selected command-line option.

\subsection{Logistic model}

The four-parameter logistic model is:
\[
y = c + \frac{L}{1 + \exp[-k(S-x_0)]}
\]

where:
\begin{itemize}[leftmargin=1.5em]
\item $L$ is the amplitude
\item $k$ controls steepness
\item $x_0$ is the transition saturation
\item $c$ is the lower or upper asymptotic baseline, depending on sign convention
\end{itemize}

In the updated comparison workflow, the logistic amplitude is allowed to be signed. This means the model can represent both:
\begin{itemize}[leftmargin=1.5em]
\item increasing relationships, where response rises with saturation
\item decreasing relationships, where response falls with saturation
\end{itemize}

This modification is particularly useful when conductivity or transformed conductivity exhibits a downward trend in the chosen analysis space.

\section{Normalization and transformation strategy}

One of the most important updates in the current framework is the handling of normalization.

If normalization is enabled, it is applied only over the data range relevant to the fitted group:
\begin{itemize}[leftmargin=1.5em]
\item sample fits use sample-specific statistics
\item site fits use site-specific statistics
\item all-data fits use the full pooled dataset
\end{itemize}

This avoids mixing scales across aggregation levels. For example, site-level fits are not built from separately normalized samples, and all-data fits are not built from separately normalized sites.

Supported normalization options are:
\begin{itemize}[leftmargin=1.5em]
\item none
\item z-score normalization
\[
x' = \frac{x-\mu}{\sigma_x}
\]
\item min--max normalization
\[
x' = \frac{x-x_{\min}}{x_{\max}-x_{\min}}
\]
\end{itemize}

When normalization is applied, both candidate models within the same group are evaluated on exactly the same transformed variables, so the AIC comparison remains internally consistent.

\section{Statistical evaluation metrics}

Model comparison is based on several complementary statistics.

\subsection{Sum of squared errors}

\[
SSE = \sum_i (y_i-\hat{y}_i)^2
\]

Smaller values indicate better agreement between observed and predicted responses.

\subsection{Coefficient of determination}

\[
R^2 = 1-\frac{SSE}{SST}
\]

where
\[
SST = \sum_i (y_i-\bar{y})^2
\]

Higher values indicate that the model explains more of the observed variance.

\subsection{Residual standard deviation}

The residual standard deviation is estimated as
\[
\sigma = \sqrt{\frac{SSE}{n}}
\]

This quantity summarizes the typical residual magnitude within a fitted group.

\subsection{Gaussian log-likelihood}

Assuming Gaussian residuals with maximum-likelihood variance estimate, the log-likelihood is
\[
\ln L =
-\frac{n}{2}
\left[
\ln(2\pi \sigma^2) + 1
\right]
\]

Larger values indicate a better fit under the Gaussian assumption.

\subsection{Akaike Information Criterion}

The Gaussian AIC used by the workflow is
\[
AIC = n\ln\left(\frac{SSE}{n}\right) + 2k
\]

where:
\begin{itemize}[leftmargin=1.5em]
\item $n$ is the number of fitted points
\item $k$ is the number of free parameters
\end{itemize}

For the power-law model, $k=2$.  
For the four-parameter logistic model, $k=4$.

The preferred model is the one with the smaller AIC.

\section{Interpretation of model choice}

The power-law and logistic models represent different physical pictures.

\subsection{Power-law interpretation}

A power-law relationship suggests smooth scaling behavior across the full saturation range. This is commonly observed for resistivity, where conduction through the fluid phase declines gradually and approximately monotonically as saturation decreases.

\subsection{Logistic interpretation}

A logistic relationship suggests a transition zone or threshold response. In conductivity data, this can reflect:
\begin{itemize}[leftmargin=1.5em]
\item onset of connected water pathways
\item enhanced pore connectivity
\item nonlinear growth of conduction after a critical saturation
\item approach toward an upper conductive plateau
\end{itemize}

Thus, model preference is not purely statistical; it also supports physical interpretation.

\section{Typical comparative patterns}

The updated workflow is especially useful for distinguishing two recurring trends:

\begin{longtable}{p{0.30\textwidth} p{0.56\textwidth}}
\toprule
\textbf{Relationship} & \textbf{Typical behavior}\\
\midrule
\endhead
Resistivity vs saturation & Often follows power-law scaling across a broad range of $S$.\\
Conductivity vs saturation & Often shows threshold-like or sigmoidal behavior, especially in real space.\\
Log-transformed resistivity & Commonly remains close to linear in log--log space.\\
Log-transformed conductivity & May either linearize or preserve evidence of nonlinear transition depending on the site and transform choice.\\
\bottomrule
\end{longtable}

The exact winner can vary by:
\begin{itemize}[leftmargin=1.5em]
\item site
\item specimen
\item normalization choice
\item transform base
\item response variable
\end{itemize}

This is why the workflow reports sample-level, site-level, and all-data results separately.

\section{Outputs from the comparison workflow}

The comparison script writes a dedicated output folder containing:
\begin{longtable}{p{0.38\textwidth} p{0.18\textwidth} p{0.31\textwidth}}
\toprule
\textbf{File} & \textbf{Scope} & \textbf{Purpose}\\
\midrule
\endhead
\texttt{stacked\_compare\_input\_\_MODE\_<tag>Hz.csv} & Campaign & Stacked data used as comparison input.\\
\texttt{compare\_power\_vs\_logistic\_\_MODE\_<tag>Hz.csv} & Campaign & One row per fitted group and response pair.\\
\texttt{compare\_power\_vs\_logistic\_\_MODE\_<tag>Hz.xlsx} & Campaign & Excel workbook with summaries and grouped tables.\\
\texttt{compare\_power\_vs\_logistic\_used\_data\_\_MODE\_<tag>Hz.csv} & Campaign & Data rows actually used in the fits.\\
PNG plots & Per case & Visual comparison of data, power-law curve, and logistic curve.\\
\bottomrule
\end{longtable}

Reported columns include:
\begin{itemize}[leftmargin=1.5em]
\item fitted parameters for both models
\item $R^2$, SSE, $\sigma$, $\ln L$, and AIC
\item model winner
\item normalization metadata
\item source folders and run names
\end{itemize}

\section{Discussion}

The expanded comparison framework makes it possible to assess model suitability more rigorously than a simple $R^2$ comparison.

A model with more parameters may achieve a better visual fit but still be penalized by AIC. This is especially relevant for logistic models in small datasets, where four parameters may be expensive unless the transition shape is truly supported by the data.

At the same time, a conductivity dataset may exhibit physically meaningful threshold behavior even when the numerical improvement is modest. For this reason, the workflow reports:
\begin{itemize}[leftmargin=1.5em]
\item goodness of fit
\item likelihood-based statistics
\item direct plotted comparisons
\end{itemize}
rather than relying on a single diagnostic.

The normalization update is also important. Without group-consistent normalization, model comparisons at site and global levels can become distorted by mixing incompatible scales from smaller subgroups.

\section{Conclusion}

This report framework supports a more complete comparison between power-law and logistic descriptions of saturation--electrical property relationships in SIP experiments.

The updated workflow contributes four major advances:
\begin{itemize}[leftmargin=1.5em]
\item comparison across sample, site, and global levels
\item optional natural-log or base-10 analysis
\item consistent normalization performed within the relevant aggregation range
\item likelihood-based statistics in addition to $R^2$
\end{itemize}

In practical terms:
\begin{itemize}[leftmargin=1.5em]
\item resistivity relationships often remain well described by power-law scaling
\item conductivity relationships often benefit from logistic representations, especially when threshold-like transitions are visible
\end{itemize}

These results reinforce the idea that different electrical properties may require different functional forms, and that model selection should be guided by both statistical evidence and physical interpretation.

\end{document}
